\documentclass[aspectratio=169,11pt]{beamer}

\usepackage[T1]{fontenc}
\usepackage[utf8]{inputenc}
\usepackage{lmodern}
\usepackage{amsmath,amssymb}
\usepackage{booktabs}
\usepackage{xcolor}
\usepackage{hyperref}
\usepackage{microtype}

\usetheme{Madrid}
\usecolortheme{default}

\definecolor{LaurierPurple}{HTML}{4B116F}
\definecolor{DeepBlue}{HTML}{0047AB}

\setbeamertemplate{navigation symbols}{}
\setbeamercolor{title}{fg=white,bg=LaurierPurple}
\setbeamercolor{frametitle}{fg=white,bg=LaurierPurple}
\setbeamercolor{block title}{fg=white,bg=DeepBlue}
\setbeamercolor{structure}{fg=DeepBlue}

\hypersetup{colorlinks=true,urlcolor=DeepBlue,linkcolor=DeepBlue}

\title[Causal ML]{Module 4: Machine Learning for Causal Inference}
\subtitle{EC 410K / EC 610I --- Machine Learning in Economics}
\author{M. Jahangir Alam, PhD}
\institute{Wilfrid Laurier University}
\date{Spring 2026}

\begin{document}

\begin{frame}[plain]
  \titlepage
\end{frame}

\begin{frame}{Module overview}
  \begin{itemize}
    \item Many causal parameters can be written so that \textbf{nuisance functions} (conditional expectations) are high-dimensional or nonlinear.
    \item \textbf{Double/debiased ML (DML)} uses ML for nuisances with \textbf{orthogonalization} and \textbf{cross-fitting} to reduce bias in estimating a low-dimensional target (e.g., \(\theta\) in a partially linear model).
    \item \textbf{Causal forests} target heterogeneous effects \(\tau(x)=\mathbb{E}[Y(1)-Y(0)\mid X=x]\) with forest-based methods.
  \end{itemize}
\end{frame}

\begin{frame}{Learning objectives}
  \begin{enumerate}
    \item State the partially linear model and interpret \(\theta\) vs the nuisance \(g(X)\).
    \item Explain why plugging in ML estimates without caution can bias causal estimates (regularization bias).
    \item Describe cross-fitting (sample splitting) at a high level.
    \item Connect methods to Chernozhukov et al.\ (2018) and Athey--Tibshirani--Wager (2019).
  \end{enumerate}
\end{frame}

\begin{frame}{Motivation: flexible confounding}
  Suppose
  \[
    Y = D\theta + g(X) + u,\qquad D = m(X) + v
  \]
  where \(g,m\) are unknown and nonlinear. A short regression of \(Y\) on \(D\) can be biased if \(D\) is correlated with \(g(X)\) through \(X\).
\end{frame}

\begin{frame}{Orthogonalization / partialling out (idea)}
  Define residuals (population):
  \[
    \tilde Y = Y - \mathbb{E}[Y\mid X],\qquad \tilde D = D - \mathbb{E}[D\mid X]
  \]
  Then (under conditional exogeneity assumptions) a regression of \(\tilde Y\) on \(\tilde D\) targets \(\theta\).
\end{frame}

\begin{frame}{Why cross-fitting?}
  \begin{itemize}
    \item If we estimate \(\widehat{\mathbb{E}}[Y\mid X]\) on the same observations we use for the second stage, overfitting can contaminate inference.
    \item \textbf{Cross-fitting:} estimate nuisances on training folds, predict on held-out folds, then second stage on orthogonalized out-of-fold predictions.
  \end{itemize}
\end{frame}

\begin{frame}{Chernozhukov et al.\ (2018)}
  \begin{itemize}
    \item Framework for \(\sqrt{n}\)-consistent estimation and inference for causal/structural parameters with ML nuisances.
    \item Key ingredients: Neyman orthogonality (insensitivity to nuisance estimation error) + cross-fitting.
  \end{itemize}
\end{frame}

\begin{frame}{Athey, Tibshirani, and Wager (2019): generalized random forests}
  \begin{itemize}
    \item Random forests adapted to estimate heterogeneous effects and other local parameters.
    \item Useful when the goal is \(\tau(x)\) for policy targeting, not only an ATE.
  \end{itemize}
\end{frame}

\begin{frame}{Replication paper connections}
  \begin{itemize}
    \item \textbf{Hansen (2015):} local heterogeneity in deterrence effects (RD setting)---conceptual link to HTE.
    \item \textbf{Anderson (2017):} effects of athletic success may vary by institution type---CATE as a substantive question.
    \item \textbf{Nunn (2008):} heterogeneity by geography/history---discuss carefully (identification remains historical).
  \end{itemize}
\end{frame}

\begin{frame}{Python / Colab demonstration}
  \begin{enumerate}
    \item Simulate nonlinear \(g(X)\), \(m(X)\) with known \(\theta\).
    \item Show naive OLS \(\mathrm{Y \sim D}\) is off.
    \item Run 2-fold cross-fitted DML with random forest nuisances; compare \(\widehat{\theta}\) to truth.
  \end{enumerate}
\end{frame}

\begin{frame}[fragile]{Colab setup}
  \begin{verbatim}
!pip -q install numpy pandas scikit-learn statsmodels matplotlib
  \end{verbatim}
\end{frame}

\begin{frame}{Student / group assignment}
  \begin{itemize}
    \item Draw the DML workflow: split \(\rightarrow\) estimate \(\widehat{\mathbb{E}}[Y\mid X]\), \(\widehat{\mathbb{E}}[D\mid X]\) \(\rightarrow\) residualize \(\rightarrow\) second stage.
    \item Present either DML (2018) or causal forests (2019) with correct scope (what assumption is still needed?).
    \item Run Colab; discuss what breaks if treatment is structurally endogenous beyond \(X\) confounding.
  \end{itemize}
\end{frame}

\begin{frame}{Summary}
  \begin{itemize}
    \item ML can estimate high-dimensional nuisances, but causal validity still requires economic identification.
    \item DML combines orthogonalization + cross-fitting to reduce regularization bias in the target parameter.
    \item Causal forests help summarize heterogeneity when assumptions for HTE interpretation hold.
  \end{itemize}
\end{frame}

\begin{frame}[plain]
  \begin{center}
    \Large Questions?\\[0.8em]
    \normalsize Next: trees, classification, interpretability (Module 5).
  \end{center}
\end{frame}

\end{document}
